\documentclass[11pt,a4paper]{article}
\usepackage[margin=23mm]{geometry}
\usepackage{amsmath,amssymb,booktabs,array}
\input{glyphtounicode}
\pdfgentounicode=1
\usepackage[hidelinks]{hyperref}
\setlength{\parindent}{0pt}
\setlength{\parskip}{6pt}
\begin{document}
\begin{center}
{\Large Step 3: information and full-time sensitivity}\\[4pt]
14 October 2026 --- analytical review under the approved H1/T1 family
\end{center}

Step 2 is user-approved and closed. This note derives review references
without changing its equations, running traffic trajectories or selecting
calibrated traffic parameters. It separates an exact model-relative
information/action possibility from externally useful advice.

\section*{1. Fixed inputs and one analytical comparison family}
Use the closed H1/T1 contract: $E$ ends, $A$ receives early transfers and
$B$ is the other continuing lane. Total road length is $L=L_p+b$, initial
footprint $[0,\ell]$, and the terminal starts empty. The common action is
$v\in[v_-,v_+]$; continuous mass is $M=\ell\rho$.
Initial current continuing densities satisfy $a_0=b_0$.
The scalar early-preparation law in physical progress is
\[
 e_z=-\frac1h\min\{e,(\delta(v)+e)/\gamma\},\quad
 \delta(v)=k_c(v)-r,\quad r=e_0+a_0,\quad
 k_c(v)=\frac{wk_j}{v+w}.
 \tag{1}
\]
The final domain has $v_+r<R\leq C(v_-)$, $v_+b_0\leq R_B\leq C(v_-)$
and $M<k_jb$. Consequently $r<k_c(v_+)$, $\delta(v)>0$, and terminal
ending sending does not bind the weighted server.

For reference speed $u>0$ and density $\kappa>0$, retain one shared geometry
and parameter family:
\[
 \begin{gathered}
 v_-=u,\quad v_+=2u,\quad w=2u,\quad k_j=3\kappa,\quad
 \gamma=2,\quad R=\frac{39}{20}\kappa u,\quad R_B=\kappa u,\\
 b=\ell,\quad L_p=\ell+h\log2,\quad L=2\ell+h\log2,\quad
 0<\eta=\ell/h<\log(4/3),\quad \alpha\geq1 .
 \end{gathered}
 \tag{2}
\]
Compare two initial preparations:
\[
 \begin{array}{c|cccc}
 &e_0&a_0=b_0&r&\rho\\ \hline
 \mathcal L&3\kappa/5&3\kappa/10&9\kappa/10&6\kappa/5\\
 \mathcal H&7\kappa/10&\kappa/4&19\kappa/20&6\kappa/5 .
 \end{array}
 \tag{3}
\]
All actions and both preparations satisfy the closed initial, ordinary-load,
exit-capacity and storage bounds: $C(v_-)=2\kappa u$ and
$M/(k_jb)=2/5$. Only initial ending mass differs as the one latent degree
under symmetry. The late work parameter is shared across preparations and
actions; no numerical $\alpha$ is selected for control experiments.

\section*{2. A reachable complete pre-decision aggregate history}
The post-decision task still has no arrivals after time zero.
To qualify initial reachability, introduce only this disclosed pre-decision
loading history, using the same H1 local balance law and a common past speed
$u$. Let $T=\ell/u$; the road is empty before $-T$.
For each preparation, during $-T\leq t\leq0$, prescribe upstream lane
density traces
\[
 e^{\rm in}_i(t)=e_{0,i}\exp(-ut/h),\quad
 a^{\rm in}_i(t)=r_i-e^{\rm in}_i(t),\quad
 b^{\rm in}_i(t)=b_{0,i}.
 \tag{4}
\]
They are exogenous preparation data, not another optimized control or future
receiving trace. The common recipe and its parameters may be given to all
predictors; the unobserved initial ending statistic selects the member.

Because $\exp(\eta)<4/3$,
\[
 e_{\mathcal H}^{\rm in}<14\kappa/15<19\kappa/20=r_{\mathcal H},
 \quad
 e_{\mathcal L}^{\rm in}<4\kappa/5<9\kappa/10=r_{\mathcal L}.
 \tag{5}
\]
All input lane densities are nonnegative and below $k_c(u)=2\kappa$.
The early donor branch is feasible throughout loading since
$a+2e=r+e<2\kappa$; the larger bound is
$19/20+14/15=113/60<2$.

At time $t$, the occupied front is $x_f=u(t+T)\leq\ell<L_p$.
Within $0<x<x_f$, the three lane densities equal the common time traces
in (4); they solve the same-coordinate H1 transport balances.
For example $e_t=-ue/h$, $e_x=0$, exactly matching the donor-limited
source. Boundary entries have those same traces, so each newly entering
characteristic joins the common homogeneous body. At time zero this
construction yields (3), the required uniform footprint and empty terminal.
Initial continuing symmetry holds at the decision instant; the loading
lane inputs are not separately asserted symmetric.

Give the aggregate observer the stronger lane-blind envelope of density,
longitudinal flow, traffic speed and cumulative counts at every cross-section
for the entire pre-decision history. In both preparations,
\[
 \begin{split}
 K_\Sigma(x,t)&=\rho\,{\bf1}_{0<x<u(t+T)},\\
 Q_\Sigma(x,t)&=uK_\Sigma(x,t),\\
 N_\Sigma(x,t)&=\rho[\,u(t+T)-x\,]_+,\qquad -T\leq t\leq0 .
 \end{split}
 \tag{6}
\]
Speed is $u$ where traffic is present; empty-channel masks are identical.
Earlier histories are empty in both cases, and all terminal counts are zero.
Thus every standard lane-blind channel above, and any function or deterministic
filter of it, is identical for all pre-decision times. A finite sensor/history
subset cannot distinguish the pair if even this stronger envelope cannot.

Lane counts, lane inputs and lateral-event counts are additional channels,
not secretly discarded longitudinal aggregate data. A lateral-event channel
could identify the preparation under this recipe. If it is actually admitted,
it must be given to A and its inference reassessed.
This is a constructed analytical reachability certificate, not evidence that
real or microsimulated data have these exact histories.

The shared pre-decision population-time prefix is
$\int_{-T}^0\rho u(t+T)\,dt=MT/2$. It is identical and independent of the
future action; excluding it from the decision-time objective changes no
comparison. No action-dependent travel or waiting is removed.

\section*{3. The compact rival exactly closes prediction}
Given $M$, $\ell$ and initial ending mass $E_0$,
\[
 e_0=E_0/\ell,\qquad
 a_0=b_0=(M-E_0)/(2\ell),\qquad r=e_0+a_0.
 \tag{7}
\]
The closed law then predicts all early densities, terminal entry/eligibility
histories, queue and departures. A+E and explicit lane representation G
therefore have identical objectives and sensitivities under this member.
Predicted histories do not constitute extra initial information.
An aggregate-only deterministic predictor receives identical inputs in
(6); an arbitrary learned prior can express uncertainty but cannot identify
both actual preparations from these channels.

\section*{4. When preparation becomes action-dependent}
For $\gamma=2$ and $\delta>0$, if $e_0\leq\delta$ the donor branch applies
for all progress and $e=e_0e^{-z/h}$. Otherwise the initial service branch
switches to the donor branch at
\[
 z_s=2h\log\frac{e_0+\delta}{2\delta},\qquad
 e(z;v)=\frac{(e_0+\delta)^2}{4\delta}e^{-z/h}
 \quad(z\geq z_s).
 \tag{8}
\]
Preparation $\mathcal L$ is donor-limited throughout every action because
$\delta_{\mathcal L}(v)\geq3\kappa/5=e_{0,\mathcal L}$.
Its completed preparation per initial position is action-independent.

For $\mathcal H$, the service branch starts when $v>v_c=18u/11$.
It ends before the first handoff exposure $z_0=L_p-\ell=h\log2$ throughout
that interval. Indeed $z_s$ is largest at $v_+$, where
$(25/22)^2<2$. The front residual $B(v)=e(v,z_0)$ is
\[
 B_{\mathcal H}(v)=
 \begin{cases}
 7\kappa/20,&v\leq v_c,\\
 (e_{0,\mathcal H}+\delta_{\mathcal H})^2/(8\delta_{\mathcal H}),
       &v>v_c .
 \end{cases}
 \tag{9}
\]
At fastest advice,
\[
 B_{\mathcal H}(v_+)=125\kappa/352,\qquad
 B'_{\mathcal H}(v_+)=225\kappa/(7744u)>0,\qquad
 B_{\mathcal L}=3\kappa/10,\quad B'_{\mathcal L}=0.
 \tag{10}
\]
The derivative in (10) is the feasible left derivative at the upper action
bound. At $v_c$, preparation's first derivative is zero on both sides.
Different residual merging alone is not yet a different speed optimum.

\section*{5. An exact full-time reference with sustained terminal waiting}
Every handoff exposure in (2)--(3) is beyond the applicable initial service
switch. Eligible lane-end arrivals begin at $t_f=(L-\ell)/v$ and last
$d=\ell/v$. Writing $s=t-t_f$ during that pulse,
\[
 \dot P_E=vB(v)e^{-vs/h},\quad
 \dot P_A=v[r-B(v)e^{-vs/h}],\quad
 \dot P_B=vb_0,\qquad 0<s<d.
 \tag{11}
\]
Define final forced mass and a pulse-area quantity
\[
 k=1-e^{-\eta},\quad m(v)=hB(v)k,\quad
 J(v)=hB(v)\ell-hm(v),\quad D_0(v)=\ell r-Rd.
 \tag{12}
\]
$m$ is mass; $J$ has units veh-eq m; $D_0$ has ordinary-work-equivalent
mass units. On the branch with positive queue throughout the eligible pulse,
\[
 q_E(s)=\frac{(vr-R)s+(\alpha-1)hB(v)(1-e^{-vs/h})}{\alpha},
 \qquad q_E(d)=\frac{D_0+(\alpha-1)m}{\alpha}.
 \tag{13}
\]
After the pulse, only ending traffic remains and it drains at $R/\alpha$.
The exact queue area and complete TST are
\[
 \begin{split}
 Q_{\rm pulse}&=\frac{D_0d}{2\alpha}
                    +\left(1-\frac1\alpha\right)\frac Jv,\\
 Q_{\rm tail}&=\frac{[D_0+(\alpha-1)m]^2}{2\alpha R},\\
 F(v)&=\frac{M(L-\ell/2)}v+Q_{\rm pulse}+Q_{\rm tail}.
 \end{split}
 \tag{14}
\]
Travel of every class and the full clearance tail are counted.
These are closed-form references, not numerically integrated traffic runs.

A simple sufficient branch condition for both preparations and all actions
in this family is $\alpha>17/3$. The donor bound gives ending density at
the latest exposure at least $9\kappa/40$ and $r\geq9\kappa/10$.
Thus the whole pulse satisfies
\[
 v[r+(\alpha-1)e]\geq
 u\kappa\left[\frac9{10}+(\alpha-1)\frac9{40}\right]>R .
 \tag{15}
\]
This branch restriction is not an empirical work-factor range.

\section*{6. A conditional information/action possibility, with an explicit limit}
Rearrange (14) exactly as
\[
 \begin{split}
 F(v)&=\frac{\alpha m(v)^2}{2R}+I_0(v)+I_1(v)/\alpha,\\
 I_0&=\frac{M(L-\ell/2)+J}{v}+\frac{(D_0-m)m}{R},\\
 I_1&=\frac{D_0d}{2}-\frac Jv+\frac{(D_0-m)^2}{2R}.
 \end{split}
 \tag{16}
\]
For $\mathcal H$ at $v_+$, $A_+=m(v_+)m'(v_+)/R>0$.
$I'_0,I'_1$ there are finite and independent of $\alpha$.
Consequently
\[
 F'_{\mathcal H}(v_+)=\alpha A_+
       +I'_0(v_+)+I'_1(v_+)/\alpha>0
 \tag{17}
\]
for the nonempty analytical parameter range
\[
 \alpha>\max\left\{\frac{17}{3},
       \frac{|I'_0(v_+)|+|I'_1(v_+)|}{A_+}\right\}.
 \tag{18}
\]
No numerical value is chosen from this range. Increasing the late work
cost eventually amplifies the extra unfinished mass at faster advice enough
to outweigh its travel advantage.

For $\mathcal L$, preparation is fixed across actions. The closed member's
earlier-arrival dominance therefore makes faster advice no worse for its
$A+E$ streams. Positive unqueued $B$ mass has strictly decreasing travel
time, so $v_+$ is its unique optimum.
For $\mathcal H$, its preparation derivative is zero near $v_-$; the same
strict travel improvement applies there. Under (18), fastest advice is not
optimal because a feasible decrease reduces $F$. Continuity on the compact
action interval therefore places its global minimizers strictly inside that
interval; uniqueness is not asserted.

The pair has identical complete aggregate histories but disjoint optimal
action sets, and opposite one-sided cost slopes at $v_+$. A common observable
policy cannot be optimal for both. The initial ending scalar resolves this
ambiguity and G adds no further predictive information.

\textbf{Limit.} This is a new conditional certificate within the approved
family, not a calibrated work-factor result, practical gradient or empirical
prevalence. The small-work result below is deliberately retained alongside it.
Intermediate work factors and other admitted geometries need their own
complete-objective classification; no positivity filter is authorized.

\subsection*{A common fastest decision for $1\leq\alpha\leq2$ in this reference family}
This negative result does not assume sustained waiting or use (14) outside
its branch. Split the local speed sensitivity into changing arrival timing
with the prepared mass profile fixed, and changing that profile at fixed
arrival times. The first part cannot increase $A+E$ time by earlier-arrival
dominance. Positive unqueued $B$ mass gives
\[
 F'_{{\rm timing},\mathcal H}\leq
 -\frac{b_{0,\mathcal H}\ell(L-\ell/2)}{v^2}
 <-\frac{\kappa h^2}{u^2}\frac{\eta}{32},
 \tag{19}
\]
using $v\leq2u$, $b_{0,\mathcal H}=\kappa/4$ and $\log2>1/2$.

For a positive queue episode, onset is the first eligible arrival because
weighted inflow decreases. At fixed clocks its preparation sensitivity is
$(\alpha-1)\partial_{v,\mathrm{prep}}P_E/\alpha$, bounded by
$(\alpha-1)m'/\alpha$. The queue episode lasts at most
$d+\alpha m/R$. A moving emptying endpoint adds no queue-area term
because its queue is zero. With no queue the local preparation contribution
is zero; the one-sided limit obeys the same bound.

In (9), $B'_{\mathcal H}$ increases from zero to its value in (10).
Indeed, putting $K_0=wk_j$, its positive-branch expression is
$(e_0^2/\delta^2-1)(\delta+r)^2/(8K_0)$; its derivative with respect to
$\delta$ is negative. Thus $m'\leq h k B'(v_+)$.
For $1\leq\alpha\leq2$, $k<\eta$, $\eta<1/3$,
$225/7744<1/32$ and $625/1716<1/2$ give
\[
 \begin{split}
 0\leq F'_{{\rm prep},\mathcal H}
 &\leq\frac12 m'_{\max}\left(\frac{\ell}{u}+\frac{2m_{\max}}R\right)\\
 &<\frac{\kappa h^2}{u^2}\frac{3\eta^2}{128}
 <\frac{\kappa h^2}{u^2}\frac{\eta}{128}.
 \end{split}
 \tag{20}
\]
Combining (19)--(20) makes complete $F_{\mathcal H}$ strictly decrease
through the action interval, including the continuous event-boundary
limits. $\mathcal L$ already has the same strict preference for every
$\alpha$. Both therefore uniquely prefer $v_+$ for $1\leq\alpha\leq2$.
This classifies the history-matched subfamily, not every H1/T1 parameter set.
Hidden preparation remains a prediction ambiguity but is not needed to
choose the optimum there.

\section*{7. The sensitivity path and later numerical evidence}
Each initial position $x$ becomes eligible at $(L-x)/v$ before terminal delay.
Let this time be $t_x(v)$; the cumulative arrival histories are
\[
 P_E(t;v)=\int_0^\ell e(v,L_p-x)
       {\bf1}_{\,t_x(v)\leq t}\,dx,\quad
 P_A(t;v)=\int_0^\ell[r-e(v,L_p-x)]
       {\bf1}_{\,t_x(v)\leq t}\,dx .
 \tag{21}
\]
$P_B$ uses constant density $b_0$. This retains both preparation sensitivity
and moving arrival times. For interior pulse times, $x_{\min}=L-vt$,
\[
 \partial_vP_E=\int_{x_{\min}}^\ell e_v(v,L_p-x)\,dx
                        +t\,e(v,vt-b).
 \tag{22}
\]
Freezing travel lags or arrival endpoints would remove the second term.
Holding initial traffic, geometry and work parameters fixed is part of
the physical control derivative.

Let $W=P_A+\alpha P_E$. The ending-queue reflection is
\[
 q_E(t;v)=\frac1\alpha\left[
 W(t;v)-Rt-\min_{0\leq s\leq t}\{W(s;v)-Rs\}\right].
 \tag{23}
\]
Eligible weighted inflow decreases during the pulse because $e$ decreases.
Therefore it admits at most one positive ending-queue excursion.
The review event inventory comprises the possible preparation branch switch,
first/last eligibility, queue onset/emptying and the pulse-end/service change.
Grazing or coincident events are not assumed smooth. Capacity and spillback
boundaries are admission boundaries, not extra clipped branches.

Choose a common fixed integration horizon larger than
$L/v_-+\alpha M/R$ for each fixed parameter case. All mass has cleared there.
Inside a differentiable event regime,
\[
 \frac{dF}{dv}=-\frac{M(L-\ell/2)}{v^2}
                   +\int_0^{H_*}\partial_v q_E(t;v)\,dt.
 \tag{24}
\]
Moving queue endpoints contribute no separate queue-area boundary term
because $q_E=0$ there, but their timing effects must already enter queue
sensitivity through (21)--(23). The final population is zero, so the
clearance endpoint does not license dropping the tail.

Later execution must first verify forward conservation/admission/clearance
and A+E/G equality. Compare the continuous analytical reference, the numerical
forward model's directional differences and its AD derivative separately.
Internal AD/difference agreement is not continuous-model convergence.
Use double precision, a stable feasible step ladder, recorded event signatures
and numerical forward/derivative convergence. At action bounds use feasible
one-sided tests. Report events separately; neither smoothing nor removal of
forward boundary cases is authorized.

Numerical tolerances, time/solver resolution and perturbation sizes must be
fixed in a later approved execution plan using the reference conditioning and
error budget. No arbitrary values or urban defaults are transplanted here.

\section*{8. External evidence and present disposition}
The proof of matched histories belongs to H1's conditional loading construction.
It is not automatically a matched-history proof for another traffic model.
The large-work range (18) has no independent empirical admission.
An independently coded use of (1), (14) or (23) is a correctness reference,
not traffic-validity or independent-usefulness evidence.

SUMO's official speed-sign, vehicle-behavior, insertion, state and trip-output
documentation was screened. It supports a potential operational speed-limit
reference with distinct vehicle interactions; it does not supply this task's
population realization, exact observer match, terminal handoff/priority,
work-factor calibration or complete evaluation mapping. A native continuous
lane-aware source model is another candidate, with different response and
service qualifications. Neither reference is adopted here.

The Step 3 source/evaluator audit and conceptual report carry the genuine
remaining choice: which independent reference meaning to qualify first.
Step 2 remains closed. Step 3 has substantive analytical findings and a
draft evidence protocol, but its full evaluation contract is not frozen.
No simulator, dependency installation, traffic experiment, AD or optimization
run was performed.
\end{document}
